\documentclass{assignment}
\usepackage{amsmath}
\usepackage{lipsum}
\usepackage{multicol}
\usepackage{fancyhdr}
\usepackage{graphicx}
\usepackage[dvipsnames]{xcolor}
\usepackage{enumitem}

\usepackage{color,soul}
\usepackage{amsfonts}
\def\code#1{\texttt{#1}}
\usepackage{graphicx} % Required for inserting images
\graphicspath{ {./images/} }
\usepackage{booktabs}
\usepackage{xcolor}
\usepackage{mdframed}
\usepackage{bbm}
\usepackage{hyperref}
\usepackage{bookmark} % Required for enhanced bookmarks
\usepackage{cleveref}

\begin{document}

\assignmentTitle{assets/kth-logo.png}
{Operating Systems}
{Assignment 3}
{ID1200}
{Lucas Frykman}
{Student Emails: lfrykman@kth.se}

\section*{Question 1}
(using p164 in OSTEP book)
\subsection*{q1(a) first-fit allocation}
\{P0, 115KB\}: S0 (32KB too small) $\rightarrow$ S1 (100KB too small) $\rightarrow$ S2 (256KB fits) $\rightarrow$ \{P0, S2\} (S2 141KB remaining)\\
\{P1, 500KB\}: All partitions are too small $\rightarrow$ \{P1, X\} \\
\{P2, 256KB\}: S0  (32KB too small) $\rightarrow$ S1 (100KB too small) $\rightarrow$ S2 (141KB too small) reaches end of list $\rightarrow$ \{P2, X\}  \\
\{P3, 31KB\}: S0 (32KB fits)  $\rightarrow$ \{P3, S0\} (S0 1KB remaining) \\ \\
Free memory remaining: (S0 1KB, S1 100KB, S2 141KB, S3 128 KB, S4 64KB)

\subsection*{q1(b) best fit allocation}
\{P0, 115KB\}: S3 (smallest partition 115KB$\le$) $\rightarrow$ \{P0, S3\} (S3 13KB remaining)\\
\{P1, 500KB\}: All partitions are too small $\rightarrow$ \{P1, X\} \\
\{P2, 256KB\}: S2 (smallest remaining partition 256KB$\le$) $\rightarrow$ \{P2, S2\} (S2 0KB remaining)\\
\{P3, 31KB\}: S0 (smallest remaining partition 32KB$\le$) $\rightarrow$ \{P3, S0\} \\ \\
Free memory remaining: (S0 1KB, S1 100KB, S2 0KB, S3 13KB, S4 64KB)

\subsection*{q1(c) worst fit allocation}
\{P0, 115KB\}: S2 (largest remaining $\ge$256KB) $\rightarrow$ \{P0, S2\} (S2 141KB remaining)\\
\{P1, 500KB\}: All partitions too small $\rightarrow$ \{P1, X\} \\
\{P2, 256KB\}: All partitions too small $\rightarrow$ \{P2, X\}  \\
\{P3, 31KB\}: S2 (largest remaining $\ge$ 31) $\rightarrow$ \{P3, S2\} (S2 110KB remaining) \\ \\
Free memory remaining: (S0 32KB, S1 100KB, S2 110KB, S3 128 KB, S4 64KB)
\subsection*{q1(d) measuring fragmentation}
External fragmentation should be a measure of how non contigious the total free memory is. When picking our method of allocation we might want to minimize this 
so that we have a higher likeliehood of succesfully allocating more processes in the near future.
In other words, of the total free memory availible we want to get as much of it as possible in to one contigious block. We can look at the largest contigious block availible and compare it 
against the rest of the total free memory.
i.e
\code{fragmentation=1-} $\frac{\code{Largest Contigious Partition}}{\code{Total Free Memory}}$ \\
Calculations for a,b,c: \\
$a = 1- \frac{141}{1+100+141+128+64} \approx 0.675$ \\

$b = 1- \frac{100}{1+100+0+13+64} \approx 0.438$ \\
$c = 1- \frac{128}{32+100+110+128+64} \approx 0.705$ \\

Best fit allocation performed the best in this example by succesfully allocating most of the processes requested as well as having the least
amount of fragmentation.

\section*{Question 2}
Assume the logical address space consists of 1024 pages with a page size of 4 KiB. The
physical memory consists of 64 frames.

\subsection*{q2(a)}
(reference p361 in textbook) \\
The logical address is a address generated by the CPU. \\
We need the $n$ bits for the page offset $d$: These bits are used to determine the page size and intra-page indexing. \\
4KiB = $4 \times 1024$ = 4096 bytes = $2^{12}$ bytes $\rightarrow n = 12$  \\
and we need $m-n$ bits for page number $p$: These bits tells us which page to look at. \\
$1024 = 2^{10} \rightarrow m-n = 10$ \\
So the logical address needs $m = 22$ bits

\subsection*{q2(b)}
To translate logical address generated by the cpu to a physical address we replace the page number $p$ with the
frame number $f$. In other words we need $\log_2(64) = 6$ $f$ bits and $n=12$ $d$ bits. 
So the physical address needs $6 + 12 = 18$ bits.

\subsection*{q2(c)}
22 bits. The virtual memry address is equivilant to the logical memory address in the context of paging. \\ \\
Now, on another system with page size of 16 KiB, calculate the page numbers and
offsets for each of the following logical address in decimal numbers \\
The decimal form of $(p,d)$ is given by \code{p*pageSize+ d}. This is simple to invert with just divison and taking the modulus. \\ \\
16KiB $= 16 \times 1024 = 16384$ bytes \\
a) Given logical address 189150 \\
Page number $p = \text{floor}(189150 / 16384) = \text{floor}(11.54) = 11$ \\
Offset number $d = 189150 \% 16384 = 8926$ \\
b) Given logical address 83205425 \\
page number $p = \text{floor}(83205425/16384) = \text{floor}(5078.45) = 5078$ \\
Offset number $d = 83205425 \% 16384 = 7473$



\section*{Question 3}

\subsection*{q3(a)}
We know $128$KiB $ = 128 \times 1024 = 131072$ bytes \\
$p_{650003} = 650003 / 131072 = \text{floor}(4.95) = 4$ \\
$d_{650003} =  650003 \% 131072 =  125715$ \\

$p_{42096} = 42096 / 131072 = 0$ \\
$d_{42096} =  42096 \% 131072 =  42096$ \\

$p_{216202} = 216202 / 131072 = \text{floor}(1.64) = 1$ \\
$d_{216202} =  216202 \% 131072 =  85130$ 
\subsection*{q3(b)}
(see p367 and lecture notes for calculating EAT) \\
For 3-level TLB we have the following hit latencies, hit ratios \\
L1 hit latency: 1ns, 0.98 \\
L2 hit latency: 1ns, 0.98 \\
L3 hit latency: 1ns, 0.98 \\
with miss latency 100ns \\
The effective access time is given by \\ EAT = \code{0.98*1ns + (0.02)*0.98*1ns + (0.02)*(0.02)*0.98*1ns + (0.02)*(0.02)*(0.02)*100ns = 1.000792ns} \\
(I'm following the convention in the lecture notes where the 100ns memory access latency isn't added to each time a hit happens.)
\subsection*{q3(c)}
(see p371) \\
We'll assume both page tables belong to a different system since the bit width of $d$ is too small for $128$KiB page size. It instead corresponds
to a page size of 4KiB = $2^{12}$ = 4096 bytes. 
We'll now construct the two-level page table using the following method of converting to binary then splitting the bits\\
\noindent\rule{10cm}{1pt} \\
$p_{650003} = 650003 / 4096 = 158 \rightarrow_{base 2}$ \code{0000 0000 0000 1001 1110} which gives us the split \\
$p_1 = (0000000000)_2 = 0$, $p_2 = (0010011110)_2 = 158$ and $d = 650003 \% 4096 = 2835$ \\
\noindent\rule{10cm}{1pt} \\
$p_{42096} = 42096 / 4096 = 10 \rightarrow_{base 2}$ \code{20b1010} which gives us the split \\
$p_1 = (0000000000)_2 = 0$, $p_2 = (0000001010)_2 = 10$ and $d = 42096 \% 4096 = 1136$ \\
\noindent\rule{10cm}{1pt} \\
$p_{216202} = 216202 / 4096 = 52 \rightarrow_{base 2}$ \code{20b1001 1110} which gives us the split \\
$p_1 = (0000000000)_2 = 0$, $p_2 = (0010011110)_2 = 52$ and $d = 216202 \% 4096 = 3210$

\subsection*{q3(d)}

\noindent\rule{10cm}{1pt} \\
$p_{650003} = 650003 / 4096 = 158 \rightarrow_{base 2}$ \code{52b1001 1110} $\rightarrow (p1,p2,p3) = (0,0,158)$
$d = 2835$ \\
\noindent\rule{10cm}{1pt} \\
$p_{42096} = 42096 / 4096 = 10 \rightarrow_{base 2}$ \code{52b1010}  $\rightarrow (p1,p2,p3) = (0,0,10)$ $d=1136$ \\
\noindent\rule{10cm}{1pt} \\
$p_{216202} = 216202 / 4096 = 52 \rightarrow_{base 2}$ \code{52b1001 1110} $\rightarrow (p1,p2,p3) = (0,0,52)$ $d = 3210$


\section*{question 4}
12 bit system with 256 byte pagesize gives page number 4 bit width and offset 8 bit width. this means that page number is
the first digit of a hex and 
\subsection*{q4(a)}
\textbf{1}. first find the page number and offset of the address \\
\code{0x2a1} = $(673)_{10} \rightarrow (p,d) = $  \\
\code{(address/pagesize,address\%pagesize)} = \code{(673/256, 673 \% 256)} $= (2, 161)$ = \code{(0x2, 0xA1)} \\
\textbf{2}. look up $p=2$ in the table and we find $f = $ \code{0xa} \\ 

\textbf{3}. construct the physical address by concatination $f \vline d =$ \code{0xaa1}  



\subsection*{q4(b)}
(\textbf{note} : I just realized at this point that the manual divison and modulus I've been doing up to this point was redundant and I could have just done all these conversions by just directly splitting up digits) \\
\textbf{1}. first find the page number and offset of the address by converting to binary then hex and split it with respect to bit widths \\
\code{0x4e6} $ \rightarrow$  \code{(0100 \vline 1110 0110)} $\rightarrow (p,d) =$ \code{(0x4,0xe6)} \\

\textbf{2}. look up $p=4$ in the table and we find that the entry is empty (i.e page fault). to handle this we must allocate a frame from the free list $\{9,f,d\}$. the first availible frame is 9. \\

\textbf{3}. construct the physical address by concatination $f \vline d =$ \code{0x9E6}  \\ 


\subsection*{q4(c)}
\textbf{1}. first find the page number and offset of the address by converting to binary then hex and split it with respect to bit widths \\
\code{0x94A} $ \rightarrow$  \code{(1001 \vline 0100 1010)} $\rightarrow (p,d) =$ \code{(0x9,0x4A)} \\

\textbf{2}. look up $p=9$ in the table and we find that $f$ = \code{0x1}  

\textbf{3}. construct the physical address by concatination $f \vline d =$ \code{0x14A}  \\ 




\section*{Question 6}
\subsection*{q6(a)}
Its a way to manage memory. Its when you only load a page when they are accessed (i.e "demanded) during the run time of a program.
Pages that are not accessed will not be loaded on to physical memory and will remain on the disk. To distinguish between
loaded/non-loaded pages you need some form of hardware support like a valid-invalid bit scheme.

\subsection*{q6(b)}

reference string: \{0,6,3,0,2,6,3,5,2,4,1,3,0,6,1,4,2,3,5,7\} \\
0  miss  [0]  fault \\ 
6 miss  [0, 6]  fault \\ 
3  miss  [0, 6, 3]  fault \\
0  hit  [0, 6, 3]  no fault \\
2  miss  [0, 6, 3, 2]  fault \\ 
6  hit  [0, 6, 3, 2]  no fault \\
3  hit  [0, 6, 3, 2]  no fault \\
5  miss  [0, 6, 3, 2, 5]  fault \\
2  hit  [0, 6, 3, 2, 5]  no fault \\ 
4  miss  [6, 3, 2, 5, 4]  evicted 0, fault \\
1  miss  [3, 2, 5, 4, 1]   evicted 6, fault \\
3  hit  [3, 2, 5, 4, 1]  no fault \\ 
0  miss  [2, 5, 4, 1, 0] 3 evicted, fault \\
6  miss  [5, 4, 1, 0, 6] evicted 2, fault \\
1  hit  [5, 4, 1, 0, 6]  No fault \\
4  hit  [5, 4, 1, 0, 6]  no fault \\ 
2  miss  [4, 1, 0, 6, 2] evicted 5, fault \\
3  miss  [1, 0, 6, 2, 3] evicted 4, fault \\
5  miss  [0, 6, 2, 3, 5] Evicted 1, fault \\
7  miss  [6, 2, 3, 5, 7] Evicted 0, fault  \\
15 page faults for first string \\

reference string: \{3,1,4,2,5,4,1,3,5,2,0,1,1,0,2,3,4,5,0,1\} \\
3 miss [3] fault \\
1 miss [3, 1] fault \\
4 miss [3, 1, 4] fault \\
2 miss [3, 1, 4, 2] fault \\
5 miss [3, 1, 4, 2, 5] fault \\
4 hit [3, 1, 4, 2, 5] no fault \\
1 hit [3, 1, 4, 2, 5] no fault \\
3 hit [3, 1, 4, 2, 5] no fault \\
5 hit [3, 1, 4, 2, 5] no fault \\
2 hit [3, 1, 4, 2, 5] no fault \\
0 miss [3, 1, 4, 2, 5] evict 3, fault \\
1 hit [1, 4, 2, 5, 0] no fault \\
1 hit [1, 4, 2, 5, 0] no fault \\
0 hit [1, 4, 2, 5, 0] no fault \\
2 hit [1, 4, 2, 5, 0] no fault \\
3 miss [4, 2, 5, 0, 3] evict 1, fault \\
4 hit [4, 2, 5, 0, 3] no fault \\
5 hit [4, 2, 5, 0, 3] no fault \\
0 hit [4, 2, 5, 0, 3] no fault \\
1 miss [2, 5, 0, 3, 1] evict 4, fault \\

8 faults

reference string \{4,2,1,7,9,8,3,5,2,6,8,1,0,7,2,4,1,3,5,8\}

4 miss [4] fault \\
2 miss [4, 2] fault \\
1 miss [4, 2, 1] failt \\
7 miss [4, 2, 1, 7] fault \\
9 miss [4, 2, 1, 7, 9] fault \\
8 miss [2, 1, 7, 9, 8] evict 4, fault \\
3 miss [1, 7, 9, 8, 3] evict 2, fault \\
5 miss [7, 9, 8, 3, 5] evict 1, fault \\
2 miss [9, 8, 3, 5, 2] evict 7, fault \\
6 miss [8, 3, 5, 2, 6] evict 9, fault \\
8 hit [8, 3, 5, 2, 6] no fault \\
1 miss [3, 5, 2, 6, 1] evict 8, fault \\
0 miss [5, 2, 6, 1, 0] evict 3, fault \\
7 miss [2, 6, 1, 0, 7] evict 5, fault \\
2 hit [2, 6, 1, 0, 7] no fault \\
4 miss [6, 1, 0, 7, 4] evict 2, fault \\
1 hit [6, 1, 0, 7, 4] no fault \\
3 miss [1, 0, 7, 4, 3] evict 6, fault \\
5 miss [0, 7, 4, 3, 5] evict 1, fault \\
8 miss [7, 4, 3, 5, 8] evict 0, fault \\
17 faults 
\end{document}
